\chapter{Accelerator}\label{ch:sw-accel}

\section{Overview}
Libre-V is intended as a host platform for testing accelerators, so the base implementation ships a
small \emph{mock accelerator} as a placeholder that a user replaces with their own design. It occupies a
subordinate endpoint on the NoC and presents a bank of 16 general-purpose 32-bit registers plus a single
interrupt line. The mock accelerator does no computation of its own; it exists to exercise, and to
document by example, the register and interrupt path an accelerator uses to talk to the cores.

The accelerator exposes a native AXI4 subordinate interface directly to the NoC (there is no APB/OBI
adapter chain, unlike the UART and timer). Its clock and reset come from a dedicated CRG domain
(\texttt{dc\_accelerator\_*}, see \autoref{ch:sw-crg}), so software must ungate the accelerator clock
and release its reset before use. Although any manager on the NoC can reach it, in the base
configuration the accelerator is driven by the Rocket core, per the platform's intended
host-CPU-plus-accelerator usage.

\begin{figure}[htbp]
  \centering
  \usetikzlibrary{positioning,arrows.meta,fit,backgrounds}
  \resizebox{\textwidth}{!}{\begin{tikzpicture}[
      font=\small,
      >={Stealth[length=2mm]},
      unit/.style={draw, rounded corners, minimum height=12mm, minimum width=26mm,
                   align=center, fill=black!5},
      lbl/.style={font=\footnotesize\ttfamily},
      node distance=14mm and 16mm
    ]
    \node[unit] (axi)  {AXI4\\ subordinate};
    \node[unit, right=of axi]  (rf)  {Register file\\ \footnotesize 16 $\times$ 32-bit};
    \node[unit, right=of rf]   (ctrl){IRQ logic\\ \footnotesize reg0 + counter};

    \coordinate (axin) at ($(axi.west)+(-16mm,0)$);
    \draw[->] (axin) -- (axi) node[midway,above,lbl]{NoC AXI};
    \draw[<->] (axi.east) -- (rf.west) node[midway,above,lbl]{r/w + strb};
    \draw[->] (rf.south) |- ($(ctrl.south)+(0,-8mm)$) -| (ctrl.south)
              node[pos=0.25,below,lbl]{reg0[1:0]};
    \coordinate (irqout) at ($(ctrl.east)+(20mm,0)$);
    \draw[->] (ctrl.east) -- (irqout) node[midway,above,lbl]{ACCEL\_INT};
  \end{tikzpicture}}
  \caption{High level schematic of the mock accelerator.}
  \label{fig:accel-schematic}
\end{figure}

\section{Register File}
The accelerator holds \texttt{NUM\_REGISTERS} (16 in the base configuration) 32-bit registers. The base
address is defined by the NoC address map (\texttt{0x3000\_0000} in Libre-V); register $i$ is at byte
offset $i \times 4$.

Because the NoC data path is 64 bits wide, the register file is addressed in \emph{even/odd pairs}: one
64-bit beat covers two consecutive 32-bit registers. The hardware forms the pair index from the address
as
\[
  \texttt{idx} = (\texttt{addr[5:2]}) \mathbin{\&} \texttt{4'b1110}
\]
so a beat carries register \texttt{idx} in its low 32 bits and register \texttt{idx+1} in its high 32
bits. A 32-bit access therefore lands in one half of a pair and is steered by the byte strobes: strobes
\texttt{[3:0]} write the even (low) register and strobes \texttt{[7:4]} write the odd (high) register.
Sub-word writes are honoured per byte, so a single-byte store modifies only its byte and leaves the
other three intact.

\begin{docnote}
Register~0 is the control register. When probing the register file with walking patterns, start at
index~1: writing bit~0 of register~0 arms the interrupt (see \autoref{sec:accel-irq}).
\end{docnote}

\section{Register Descriptions}
The registers share one layout: register~0 has the defined control bits below, and registers 1--15 are
plain 32-bit read/write storage.

\regtitle{0}{Accelerator Control Register (CTRL)}

\regmeta{0x3000\_0000}{0x000}{0x00000000}{RW}{Only bits [1:0] are functional; bits [31:2] are storage.}

\begin{register}{H}{Accelerator Control Register (CTRL)}{}

  \regfield[gray!20]{Reserved}{30}{2}{{0}}
  \regfield{IRQ\_CLR}{1}{1}{0}
  \regfield{IRQ\_EN}{1}{0}{0}
  \reglabel{Reset}

  \regfieldvalue[gray!20]{Reserved}{30}{2}{{RW}}
  \regfieldvalue{IRQ\_CLR}{1}{1}{{RW}}
  \regfieldvalue{IRQ\_EN}{1}{0}{{RW}}
  \reglabel{Type}

\end{register}

\textbf{Functional Description.}

Register~0 arms and clears the accelerator interrupt. \textbf{IRQ\_EN} (bit~0) arms the interrupt engine;
\textbf{IRQ\_CLR} (bit~1) acknowledges an asserted interrupt. The upper 30 bits are stored and read back
but have no effect. The behaviour of the interrupt engine is described in \autoref{sec:accel-irq}.

\begin{regmap}
31:2 & reserved & RW & 0 &
Reserved. Stored but have no effect. \\
1    & IRQ\_CLR  & RW & 0 &
Interrupt clear/acknowledge. Writing 1 deasserts an active interrupt and restarts the counter; the bit
self-clears in hardware. \\
0    & IRQ\_EN   & RW & 0 &
Interrupt enable (arm). While 1, the engine counts and asserts the interrupt; while 0, the interrupt is
held deasserted and the counter is reset. \\
\end{regmap}

\regtitle{$i$}{General-Purpose Register (REG$i$), $i = 1 \ldots 15$}

\regmeta{0x3000\_0000}{$i \times$ 0x004}{0x00000000}{RW}{Plain 32-bit storage. Byte-strobe writable.}

\begin{register}{H}{General-Purpose Register (REG$i$)}{}

  \regfield{DATA}{32}{0}{{00000000000000000000000000000000}}
  \reglabel{Reset}

  \regfieldvalue{DATA}{32}{0}{{RW}}
  \reglabel{Type}

\end{register}

\textbf{Functional Description.}

Registers 1 through 15 are general-purpose 32-bit scratch registers with no hardware side effects. In
the mock accelerator they demonstrate the read/write and byte-strobe paths; in a real accelerator they
are where the design places its command, status, and data registers.

\begin{regmap}
31:0 & DATA & RW & 0 &
General-purpose data. Full 32-bit read/write, byte and half-word strobes honoured. \\
\end{regmap}

\section{Interrupt}\label{sec:accel-irq}
The accelerator drives a single interrupt line, delivered to the Rocket's PLIC as source~2
(\texttt{ROCKET\_PLIC\_SRC\_ACCEL}, see \autoref{ch:sw-rocket-sys}). The engine is a simple counter
whose limit is \texttt{ACCEL\_INT\_COUNT} (100 cycles in the base configuration):

\begin{itemize}[nosep]
  \item With \textbf{IRQ\_EN} = 0 the interrupt is held low and the counter is held at 0.
  \item When \textbf{IRQ\_EN} is set to 1 the counter advances each cycle. After
        \texttt{ACCEL\_INT\_COUNT} cycles the interrupt line asserts and \emph{stays} asserted (it is a
        level, not a pulse).
  \item Writing \textbf{IRQ\_CLR} = 1 acknowledges the interrupt: the line deasserts, \textbf{IRQ\_CLR}
        self-clears, and the counter restarts. While the accelerator remains armed it therefore
        re-asserts periodically, once every \texttt{ACCEL\_INT\_COUNT} cycles.
\end{itemize}

\begin{docwarning}
Because the interrupt is a level that re-arms while \textbf{IRQ\_EN} stays 1, an in-flight assertion can
re-appear if the device is cleared without first masking it at the PLIC. To disarm cleanly, mask the
source at the PLIC \emph{first}, then clear \textbf{IRQ\_EN} at the device. A handler should acknowledge
with \textbf{IRQ\_CLR} alone (bit~1), not \texttt{IRQ\_EN|IRQ\_CLR}, so it does not re-arm.
\end{docwarning}

\section{Address Map Summary}

\begin{tabular}{@{}lll@{}}
\toprule
\textbf{Offset} & \textbf{Register} & \textbf{Access} \\
\midrule
0x000        & CTRL (register 0: IRQ\_EN, IRQ\_CLR) & RW \\
0x004--0x03C & REG1--REG15 (general purpose)        & RW \\
\bottomrule
\end{tabular}

\section{Source Files}

\begin{tabular}{@{}ll@{}}
\toprule
\textbf{File} & \textbf{Role} \\
\midrule
\texttt{accelerator/src/accelerator.sv} & Mock accelerator: register file, byte strobes, IRQ counter \\
\texttt{top/noc\_accelerator.sv}        & NoC-side wrapper (struct-to-flat AXI4 subordinate) \\
\bottomrule
\end{tabular}

\newpage
